\documentclass[11pt,a4paper]{article}
\usepackage{gtu-el-yazisi}
\begin{document}
\notbilgi{Görüntü İşleme}{Örnekler ve Image Segmentation (kenar, gradyan, eşikleme)}{Ders7.pdf}

\sayfa{1}
\[ g(x,y)=f(x,y)*\underbrace{h_1(x,y)*h_2(x,y)}_{h(x,y)}-f(x,y) \]
\[ G(u,v)=F(u,v)\left[H(u,v)-1\right] \]
\elcizimi[0.25]{sekiller/Ders7-s1-1.png}[1]
\elcizimi[0.43]{sekiller/Ders7-s1-2.png}[1]
\elcizimi[0.4]{sekiller/Ders7-s1-3.png}[1]
\elcizimi[0.6]{sekiller/Ders7-s1-4.png}[1]

\sayfa{2}
\elcizimi[0.25]{sekiller/Ders7-s2-1.png}[2]
\begin{center}
\renewcommand{\arraystretch}{1.25}
\begin{tabular}{c|c}
$z$ & $P_z$ \\ \hline
0 & 0 \\
1 & 2 \\
2 & 2 \\
3 & 2 \\
4 & 1 \\
5 & 1 \\
6 & 1 \\
7 & 0
\end{tabular}
\qquad
\begin{tabular}{c|c|c|c l}
$r$ & $p_r(r)$ & $s=T_r(r)$ & $G(z)$ & \\ \hline
0 & 1/9 & 1/9 & 0/9 & \textcolor{kirmizikalem}{0,1 $\rightarrow$ {\scriptsize tutarlı almak için}}\\
1 & 3/9 & 4/9 & 2/9 & \textcolor{kirmizikalem}{2}\\
2 &  & 6/9 & 4/9 & \textcolor{kirmizikalem}{3}\\
3 &  & 7/9 & 6/9 & \textcolor{kirmizikalem}{4}\\
4 &  & 7/9 & 7/9 & \textcolor{kirmizikalem}{4}\\
5 &  & 7/9 & 8/9 & \textcolor{kirmizikalem}{4}\\
6 &  & 7/9 & 9/9 & \textcolor{kirmizikalem}{4}\\
7 &  & 9/9 & 9/9 & \textcolor{kirmizikalem}{7}
\end{tabular}
\end{center}
3. soru \quad $(-1)^{x+y}$

\elcizimi[0.85]{sekiller/Ders7-s2-2.png}[2]
\elcizimi[0.35]{sekiller/Ders7-s2-3.png}[2]
\[ f(x,y)=\sum_{u}\sum_{v}F(u,v)\,e^{j2\pi\left(\frac{ux}{M}+\frac{vy}{N}\right)} \]

\sayfa{3}
\section{Image Segmentation}
\sekilnotu[3]{Bu sayfa basılı materyal içeriyor (Bölüm 10.1 Fundamentals, Şekil 10.1, kitaptan); dijitalleştirilmedi.}

(a) $\rightarrow$ Tüm bölgelerin toplamı image'ı oluşturmalı

(c) $\rightarrow$ Her bölgenin kendi pikselleri var

(e) $\uparrow$ \okunamadi{her yer} edge

(f) $\uparrow$ varyanstan bulunabilir $3\times3$ pencereyle

\sayfa{4}
\sekilnotu[4]{Bu sayfa basılı materyal içeriyor (Şekil 10.7 ve 10.8, kitaptan); dijitalleştirilmedi.}
\elcizimi[0.3]{sekiller/Ders7-s4-1.png}[4]
double edge sonuç var

\noindent Sınırları belirlemek, bölgeyi bulmak için edge detection\\
\hspace*{2em}segmentation \hfill gerekiyor

(b) $\uparrow$ Merkezi bulabiliyor olman lazım

\sayfa{5}
\sekilnotu[5]{Bu sayfa basılı materyal içeriyor (Şekil 10.10 ve 10.11, kitaptan); dijitalleştirilmedi.}

Zero crossing $\rightarrow$ pozitiften negatife geçişin olduğu yer

\elcizimi[0.38]{sekiller/Ders7-s5-1.png}[5]
birinci dereceden türevinin mutlağı 1'den büyük olmalı

\sayfa{6}
\[
\nabla f=\begin{bmatrix} g_x\\ g_y\end{bmatrix}=\begin{bmatrix}\dfrac{\partial f}{\partial x}\\[6pt] \dfrac{\partial f}{\partial y}\end{bmatrix}
\]
\begin{align*}
M(x,y) &= \mathrm{mag}(\nabla f)=\left(g_x^{2}+g_y^{2}\right)^{1/2}\\
&\approx |g_x|+|g_y|
\end{align*}
direction of gradient vector
\[ \alpha(x,y)=\tan^{-1}\!\left(\frac{g_y}{g_x}\right) \]
\elcizimi[0.38]{sekiller/Ders7-s6-1.png}[6]
\sekilnotu[6]{Bu sayfa basılı materyal içeriyor (Şekil 10.16 ve 10.36, kitaptan); dijitalleştirilmedi.}

$\nearrow$ iki ayrı görüntü var \okunamadi{diğerinde}

\sayfa{7}
\sekilnotu[7]{Bu sayfa basılı materyal içeriyor (Şekil 10.37 ve 10.38, kitaptan); dijitalleştirilmedi.}
\section{Basic Intensity Tresholding}
\[
g(x,y)=\begin{cases}1 & \text{if } f(x,y)>T\\ 0 & \text{o.w}\end{cases}
\]
Global tek eşik

\end{document}
